% =====================================================================
% TCC – Empresas familiares vs. não familiares: comportamentos de endividamento
% Renan do Amaral Ribeiro – FEA-RP/USP – Ciências Contábeis
% Classe abnTeX2 (ABNT NBR 14724, 6023, 10520, 6027, 6028)
% Compilar: latexmk  (ou: make)   →  gera build/main.pdf
% =====================================================================
\documentclass[
  12pt,
  oneside,          % impressão só frente (use twoside se for frente e verso)
  a4paper,
  english,          % idioma adicional (abstract)
  brazil            % idioma principal: o último da lista
]{abntex2}

% ---------------------------------------------------------------------
% Pacotes
% ---------------------------------------------------------------------
\usepackage[T1]{fontenc}
\usepackage[utf8]{inputenc}
\usepackage{lmodern}
\usepackage{indentfirst}
\usepackage{microtype}
\usepackage{xcolor}
\usepackage{graphicx}
\usepackage{booktabs}
\usepackage{tabularx}
\usepackage{array}
\usepackage{enumitem}
\usepackage{amsmath}
\usepackage[brazilian,hyperpageref]{backref}
\usepackage[alf]{abntex2cite}   % citações autor-data (NBR 10520)

% ---------------------------------------------------------------------
% Marcadores de revisão (aparecem em vermelho no PDF)
%   Para a versão final, troque \revisaotrue por \revisaofalse:
%   todos os marcadores somem do PDF.
% ---------------------------------------------------------------------
\newif\ifrevisao
\revisaotrue
\ifrevisao
  \newcommand{\verificar}[1]{\textcolor{red}{\textbf{[VERIFICAR:} #1\textbf{]}}}
  \newcommand{\decidir}[1]{\textcolor{orange!80!black}{\textbf{[DECIDIR:} #1\textbf{]}}}
\else
  \newcommand{\verificar}[1]{}
  \newcommand{\decidir}[1]{}
\fi

% Ambiente "quadro" e Lista de Quadros (receita da documentação do abnTeX2)
\newcommand{\quadroname}{Quadro}
\newcommand{\listofquadrosname}{Lista de quadros}
\newfloat[chapter]{quadro}{loq}{\quadroname}
\newlistof{listofquadros}{loq}{\listofquadrosname}
\newlistentry{quadro}{loq}{0}
\counterwithout{quadro}{chapter}
\renewcommand{\cftquadroname}{\quadroname\space}
\renewcommand*{\cftquadroaftersnum}{\hfill--\hfill}
\setfloatlocations{quadro}{hbtp}

% Colunas de tabela com quebra de linha alinhada à esquerda
\newcolumntype{L}{>{\raggedright\arraybackslash}X}
\newcolumntype{P}[1]{>{\raggedright\arraybackslash}p{#1}}

% ---------------------------------------------------------------------
% Dados do trabalho
% ---------------------------------------------------------------------
\titulo{Empresas familiares vs. não familiares: comportamentos de endividamento}
\autor{Renan do Amaral Ribeiro}
\local{Ribeirão Preto}
\data{2026}
\orientador[Orientadora:]{Profa. Dra. Maisa Ribeiro}
\instituicao{%
  Universidade de São Paulo
  \par
  Faculdade de Economia, Administração e Contabilidade de Ribeirão Preto
  \par
  Departamento de Contabilidade}
\tipotrabalho{Trabalho de Conclusão de Curso (Graduação)}
\preambulo{Trabalho de Conclusão de Curso apresentado ao Departamento de
Contabilidade da Faculdade de Economia, Administração e Contabilidade de
Ribeirão Preto da Universidade de São Paulo, como requisito parcial para a
obtenção do título de Bacharel em Ciências Contábeis.}

% PDF
\makeatletter
\hypersetup{
  pdftitle={\@title},
  pdfauthor={\@author},
  pdfsubject={Trabalho de Conclusão de Curso},
  pdfkeywords={empresas familiares}{endividamento}{estrutura de capital},
  colorlinks=true,
  linkcolor=black,
  citecolor=black,
  urlcolor=blue!50!black,
  bookmarksdepth=4
}
\makeatother

% Espaçamento de parágrafo e recuo
\setlength{\parindent}{1.3cm}
\setlength{\parskip}{0.2cm}

% Lista de quadros e texto da referência cruzada das citações
\renewcommand{\backrefpagesname}{Citado na(s) página(s):~}
\renewcommand{\backref}{}
% "Citado na página X" nas referências: só na versão de revisão (não é exigido pela ABNT)
\ifrevisao
  \renewcommand*{\backrefalt}[4]{%
    \ifcase #1 Nenhuma citação no texto.%
    \or Citado na página #2.%
    \else Citado #1 vezes nas páginas #2.%
    \fi}
\else
  \renewcommand*{\backrefalt}[4]{}
\fi

% Capa no modelo ABNT (NBR 14724): instituição, autor, título, local e ano
\renewcommand{\imprimircapa}{%
  \begin{capa}%
    \center
    {\ABNTEXchapterfont\large UNIVERSIDADE DE SÃO PAULO\par
    FACULDADE DE ECONOMIA, ADMINISTRAÇÃO E CONTABILIDADE\par
    DE RIBEIRÃO PRETO\par
    DEPARTAMENTO DE CONTABILIDADE\par}
    \vspace*{3cm}

    {\ABNTEXchapterfont\large\imprimirautor}

    \vfill
    \begin{center}
    \ABNTEXchapterfont\bfseries\LARGE\imprimirtitulo
    \end{center}
    \vfill

    \large\imprimirlocal

    \large\imprimirdata

    \vspace*{1cm}
  \end{capa}
}

\makeindex

% =====================================================================
\begin{document}
\selectlanguage{brazil}
\frenchspacing

% ---------------------------------------------------------------------
% Elementos pré-textuais
% ---------------------------------------------------------------------
\pretextual
\imprimircapa
\imprimirfolhaderosto
\cleardoublepage % necessário: sem a quebra, o ambiente resumo falha nesta versão do abnTeX2
% \include{pretextuais/ficha}       % ficha catalográfica, se exigida
% \include{pretextuais/aprovacao}   % folha de aprovação, se exigida
% Resumo – NBR 6028:2021 (parágrafo único, 150 a 500 palavras, voz ativa, 3.ª pessoa)
\begin{resumo}
Este trabalho investiga se empresas familiares brasileiras de capital aberto
apresentam endividamento diferente do das empresas não familiares. Com base na
perspectiva da riqueza socioemocional, a pesquisa parte da hipótese de que a
aversão à interferência de credores e à perda de controle leva as famílias
controladoras a adotarem políticas de financiamento mais conservadoras. A amostra
compreende 16 companhias do setor de saúde da B3, com dados trimestrais da
Economática de 2018 a 2025. As empresas foram classificadas como familiares
segundo o critério de \citeonline{santos2018}, e dois critérios alternativos foram
usados como testes de robustez. O endividamento foi medido por indicadores
contábeis de nível, de capacidade de pagamento e de composição da dívida, e as
diferenças entre os grupos foram avaliadas pelos testes t de Welch e de
Mann-Whitney, para o período completo e para os subperíodos anterior, durante e
posterior à pandemia de COVID-19. Os resultados indicam que as empresas
familiares apresentaram, em todos os exercícios, menor endividamento oneroso e
menor participação da dívida de curto prazo que as não familiares, mas essas
diferenças não foram estatisticamente significativas no período completo, e o
endividamento total dos dois grupos foi semelhante. A exceção é a composição da
dívida durante a pandemia, quando a menor participação da dívida de curto prazo
nas empresas familiares foi significativa ao nível de 5\%, em análise exploratória
por subperíodos, sem correção para múltiplos testes. A diferença no endividamento
oneroso inverteu o sinal sob os critérios alternativos de classificação e
mostrou-se sensível à exclusão de uma empresa sem dívida financeira. Conclui-se que
a hipótese não pode ser confirmada estatisticamente. A evidência descritiva é
compatível com uma postura mais conservadora das famílias em relação ao prazo da
dívida, mas também com diferenças de tamanho e de acesso ao crédito entre os
grupos.

\textbf{Palavras-chave}: empresas familiares; endividamento; estrutura de capital;
riqueza socioemocional; setor de saúde.
\end{resumo}

% Abstract – resumo em língua estrangeira (NBR 14724).
% RASCUNHO escrito a partir do resumo em português, não traduzido palavra por palavra.
% Revise antes da entrega (de preferência com alguém fluente em inglês).
\begin{resumo}[Abstract]
\begin{otherlanguage*}{english}
This study examines whether Brazilian listed family firms differ from non-family
firms in their use of debt. Drawing on the socioemotional wealth perspective, it
hypothesizes that controlling families, averse to creditor interference and to
the loss of control, adopt more conservative financing policies. The sample
comprises 16 healthcare companies listed on B3, with quarterly Economática data
from 2018 to 2025. Firms were classified as family firms following
\citeonline{santos2018}, and two alternative definitions were used as robustness
checks. Debt was measured by accounting ratios of leverage, debt service capacity
and debt maturity, and group differences were assessed with Welch's t test and the
Mann-Whitney test over the full period and over the pre-pandemic, pandemic and
post-pandemic subperiods. Family firms showed lower financial debt and a smaller
share of short-term debt in every year, but these differences were not
statistically significant over the full period, and total liabilities were
similar in both groups. The exception is debt maturity during the COVID-19
pandemic, when the smaller share of short-term debt in family firms was
significant at the 5\% level, in an exploratory subperiod analysis without
correction for multiple testing. The difference in financial debt changed sign
under the alternative family firm definitions and was sensitive to the exclusion
of one firm with no financial debt. The hypothesis is therefore not statistically
supported. The descriptive evidence is consistent with a more conservative stance
of families towards debt maturity, but also with differences in firm size and
access to credit between the groups.

\vspace{\onelineskip}
\noindent
\textbf{Keywords}: family firms; debt; capital structure; socioemotional wealth;
healthcare sector.
\end{otherlanguage*}
\end{resumo}


% Listas: descomente a de ilustrações e a de tabelas quando houver figuras e tabelas
% (uma lista vazia não deve aparecer no trabalho).
\pdfbookmark[0]{\listfigurename}{lof}
\listoffigures*
\cleardoublepage
\pdfbookmark[0]{\listofquadrosname}{loq}
\listofquadros*
\cleardoublepage
\pdfbookmark[0]{\listtablename}{lot}
\listoftables*
\cleardoublepage
\pdfbookmark[0]{\contentsname}{toc}
\tableofcontents*
\cleardoublepage

% ---------------------------------------------------------------------
% Elementos textuais
% ---------------------------------------------------------------------
\textual
% =====================================================================
\chapter{Introdução}\label{cap:introducao}
% =====================================================================

O ambiente corporativo brasileiro é composto por empresas com diferentes
estruturas de propriedade e de governança. Entre elas, destacam-se as empresas
familiares, que representam uma parcela relevante das organizações e possuem
características particulares relacionadas à preservação do controle, sucessão e
continuidade do negócio ao longo das gerações \cite{pestana2021}. Em contraste,
empresas com estruturas acionárias mais dispersas tendem a apresentar modelos de
gestão mais profissionalizados e menor influência de um grupo controlador
específico. Neste cenário, as decisões de estrutura de capital, isto é, a
combinação entre capital próprio e de terceiros, representam um dilema central da
gestão financeira. Teorias clássicas como o \textit{Trade-Off}, que explica a
escolha ótima de endividamento com base no equilíbrio entre benefícios fiscais e
custos de falência, e a \textit{Pecking Order}, que propõe uma hierarquia de
financiamento baseada na preferência por recursos internos, oferecem
interpretações racionais para essas decisões, assumindo como objetivo primordial a
maximização do valor da firma \cite{myers1984,myersmajluf1984}.

Contudo, a aplicação direta destes modelos torna-se complexa no contexto de
empresas familiares. A literatura sobre governança em firmas familiares introduz o
conceito de riqueza socioemocional (\textit{socioemotional wealth} -- SEW), que
representa a utilidade não financeira que a família extrai do negócio
\cite{gomezmejia2007}. Essa
``riqueza'' abrange a manutenção do controle, a perpetuidade do legado e a
identidade familiar atrelada à organização, objetivos que podem se sobrepor à pura
maximização de resultados econômicos \cite{santos2018}.

Essa busca pela preservação do SEW colide diretamente com os riscos inerentes ao
endividamento. Contratos de dívida, especialmente os de longo prazo, frequentemente
impõem cláusulas restritivas (\textit{covenants}) que servem como mecanismo de
disciplina, mas que também representam uma potencial ameaça à autonomia gerencial
\cite{costa2018}. A violação dessas cláusulas pode resultar em penalidades como a
renegociação compulsória das obrigações financeiras ou até mesmo o vencimento
antecipado da dívida. Nesse contexto, verificou-se que empresas com níveis mais
elevados de endividamento podem apresentar incentivos ao gerenciamento de
resultados, buscando reduzir a probabilidade de descumprimento dos
\textit{covenants} e os custos associados à sua violação \cite{costa2018}. Para
uma família controladora, a mera possibilidade de intervenção de credores
representa uma ameaça direta ao seu poder decisório.

Assim, a aversão à interferência de credores, que ameaça o SEW, pode levar a uma
política de endividamento mais conservadora. Como o
setor de atuação também influencia a estrutura de capital \cite{oliveira2020}, a
comparação entre empresas familiares e não familiares precisa reduzir o risco de
que diferenças setoriais sejam confundidas com o efeito da estrutura de controle. Por
isso, esta pesquisa concentra a análise nas companhias de um único setor, o de
saúde da B3 S.A. -- Brasil, Bolsa, Balcão (B3).
\verificar{conferir em Oliveira e Kayo (2020) a passagem sobre o efeito do setor}

Diante do exposto, emerge o seguinte problema de pesquisa: a estrutura de controle
familiar está associada a diferentes níveis de endividamento quando comparada à
estrutura de controle não familiar nas companhias abertas brasileiras do setor de
saúde? A resposta a esta questão pode subsidiar investidores, credores e os
próprios gestores familiares na formulação de estratégias financeiras alinhadas a
seus objetivos econômicos e não econômicos.

O objetivo geral deste trabalho é comparar o nível e a composição do endividamento
de empresas familiares e não familiares do setor de saúde da B3 no período de 2018
a 2025. Como objetivos específicos, busca-se: (i) classificar as companhias quanto
ao controle familiar a partir de critérios da literatura; (ii) comparar os grupos
quanto ao endividamento oneroso, ao endividamento total, à capacidade de pagamento
e, de forma complementar e exploratória, ao prazo da dívida; e (iii) verificar se as diferenças se mantêm em subperíodos e
sob critérios alternativos de classificação.

Nesse contexto, justifica-se a realização deste estudo pela relevância prática e
teórica do tema. Apesar da extensa literatura sobre estrutura de capital e
empresas familiares, representada por estudos como \citeonline{kayo2018},
\citeonline{wildner2020} e, em empresas portuguesas, \citeonline{pestana2021}, os
resultados acerca da
influência da estrutura de controle familiar sobre o endividamento corporativo
permanecem inconclusivos. Enquanto alguns estudos apontam para uma postura mais
conservadora e menores níveis de alavancagem em empresas familiares, outros
identificam comportamentos distintos em função de características específicas de
governança, do ambiente institucional e das particularidades de cada mercado. Além
disso, esta pesquisa examina um período recente, que inclui o choque da pandemia de
COVID-19, concentra a comparação em um único setor e verifica se os resultados se
mantêm sob diferentes critérios de classificação de empresa familiar, aspecto em
que os estudos nacionais divergem. Com isso, busca oferecer uma contribuição
empírica adicional para a literatura de finanças corporativas e de governança em
empresas familiares. Do ponto de vista gerencial, os resultados podem auxiliar
gestores, credores e investidores a compreender de forma mais clara como a
presença da família no controle afeta o apetite ao endividamento e,
consequentemente, o risco financeiro da organização.

Além desta introdução, o trabalho está organizado da seguinte forma: o
\autoref{cap:referencial} apresenta o referencial teórico e a hipótese de
pesquisa; o \autoref{cap:metodologia} descreve a amostra, a classificação das
empresas, os indicadores e os procedimentos estatísticos; o
\autoref{cap:resultados} apresenta e discute os resultados; e o
\autoref{cap:consideracoes} traz as considerações finais.

% =====================================================================
\chapter{Referencial teórico}\label{cap:referencial}
% =====================================================================

O referencial teórico a seguir constrói a base conceitual que suporta a hipótese
de menor endividamento em empresas familiares, partindo das teorias clássicas e
aprofundando-se nas particularidades de sua governança.

% ---------------------------------------------------------------------
\section{Conflitos de agência e a riqueza socioemocional (SEW)}\label{sec:sew}
% ---------------------------------------------------------------------

O arcabouço teórico que define as empresas familiares enfatiza o envolvimento da
família no controle, propriedade e gestão \cite{santos2018}. Entre as companhias
analisadas nesta pesquisa, a Baumer ilustra a presença simultânea dessas três
dimensões: a família detém cerca de 89\% das ações ordinárias, integra o bloco de
controle e ocupa a presidência do conselho de administração e a
diretoria-presidência. Já a RaiaDrogasil exemplifica um arranjo comum no mercado
brasileiro, em que três famílias reunidas em um mesmo bloco de controle detêm, em
conjunto, menos da metade das ações ordinárias, mas mantêm membros no conselho de
administração.
\verificar{participação e cargos atuais na BALM e na RADL no Formulário de
Referência mais recente}

Em contraste com o modelo de agência clássico, que foca nos conflitos entre
gestores e acionistas, estudos sobre empresas familiares destacam que essas
organizações frequentemente perseguem objetivos que transcendem a maximização do
retorno econômico, incorporando aspectos relacionados à preservação do controle,
da identidade familiar e da continuidade do negócio ao longo das gerações
\cite{ames2020,santos2018}. Nesse contexto, ganha destaque o conceito de riqueza
socioemocional (SEW), que abrange elementos como a manutenção da identidade
familiar, a autonomia decisória e o desejo de perpetuidade do negócio
\cite{santos2018}.

\citeonline{santos2018}, ao analisarem a remuneração executiva, fornecem
evidências de que a prioridade em preservar o SEW pode se manifestar em incentivos
menos intensos em remuneração pecuniária, com remuneração total e variável
menores, sugerindo um alinhamento de interesses intrínseco que minimiza o conflito
clássico entre gestor e proprietário. Contudo, a prevalência do SEW gera um
conflito latente com os credores, uma vez que a dívida, por meio de mecanismos de
monitoramento e restrições contratuais, pode representar uma ameaça à manutenção
do controle e da autonomia da família controladora, aspectos considerados centrais
para a preservação da riqueza socioemocional \cite{santos2018,ames2020}.

% ---------------------------------------------------------------------
\section{Endividamento elevado e o risco de interferência}\label{sec:risco}
% ---------------------------------------------------------------------

A dívida corporativa, embora ofereça benefícios fiscais, eleva o risco de
dificuldade financeira. Um dos mecanismos contratuais mais relevantes nesse
contexto é o \textit{covenant}, uma cláusula restritiva presente em contratos de
dívida que estabelece limites para determinados indicadores
econômico-financeiros e atua como instrumento de monitoramento e proteção dos
credores \cite{costa2018}.

A relação entre endividamento e escolhas contábeis não é linear.
\citeonline{costa2018} encontraram evidências de que, ao atingir níveis muito
elevados de endividamento, as empresas brasileiras tendem a aumentar o componente
discricionário dos lucros para evitar a quebra de cláusulas contratuais. Nesse
contexto, a redução da qualidade dos lucros pode ser percebida pelos gestores como
menos onerosa do que os custos decorrentes da violação dos \textit{covenants},
incluindo a renegociação compulsória ou o vencimento antecipado das obrigações
financeiras. Essa relação de risco é especialmente associada às dívidas de longo
prazo, que normalmente apresentam cláusulas restritivas mais rigorosas.

O prazo da dívida, contudo, expõe a empresa aos credores de outra forma. Enquanto a
dívida de longo prazo concentra as restrições nas cláusulas contratuais, a dívida
de curto prazo precisa ser renovada com frequência e sujeita a empresa ao risco de
liquidez, isto é, ao risco de que os credores se recusem a refinanciá-la ou o
façam em condições desfavoráveis \cite{diamond1991}. Esse risco tende a se agravar
em períodos de crise, quando o crédito se torna mais escasso. Para uma família que
busca preservar a autonomia decisória, depender de renovações frequentes significa
submeter-se repetidamente à avaliação dos credores. Assim, além de um menor nível
de dívida, a aversão à interferência de credores pode se manifestar na preferência
por prazos mais longos.

Para as empresas familiares, o risco de violação de \textit{covenants} pode ser
particularmente relevante, uma vez que a eventual intervenção dos credores pode
limitar a autonomia decisória da família controladora. Considerando que a
preservação do controle e da independência da gestão constitui um dos principais
componentes da riqueza socioemocional, a exposição excessiva ao endividamento pode
representar uma ameaça a objetivos que transcendem o desempenho econômico da firma
\cite{santos2018,pestana2021}. Portanto, espera-se que a família controladora
apresente maior aversão à alavancagem excessiva como forma de reduzir o risco de
perda de controle.

Há, porém, um argumento em sentido oposto. Como a emissão de novas ações dilui a
participação da família, a preservação do controle pode levar a família a preferir
a dívida ao capital próprio. \citeonline{poletti2022}, em empresas
latino-americanas, encontraram empresas familiares mais alavancadas quando a
participação da família é menor, como forma de manter o controle, e redução da
alavancagem quando essa participação é elevada. No Brasil, \citeonline{kayo2018}
observaram empresas familiares mais endividadas e com ajuste mais lento da
estrutura de capital. Essas evidências ajudam a explicar por que a literatura
permanece inconclusiva. Ainda assim, considerando a perspectiva da riqueza
socioemocional e o papel dos credores discutido nesta seção, adota-se como
hipótese a direção de menor endividamento.
\verificar{conferir nos artigos de Poletti-Hughes e Martínez Garcia (2022) e de
Kayo, Brunaldi e Aldrighi (2018) os resultados descritos neste parágrafo}

Com base no exposto, formula-se a seguinte hipótese:

\begin{quote}
\textbf{H1: Empresas de controle familiar apresentam um nível de alavancagem
financeira inferior ao de empresas não familiares.}
\end{quote}

% =====================================================================
\chapter{Metodologia}\label{cap:metodologia}
% =====================================================================

Este capítulo descreve o delineamento da pesquisa, a amostra e as fontes de dados,
o critério de classificação das empresas familiares, os indicadores de
endividamento e os procedimentos estatísticos empregados para avaliar a hipótese
H1.

% ---------------------------------------------------------------------
\section{Abordagem e estratégia da pesquisa}\label{sec:abordagem}
% ---------------------------------------------------------------------

A pesquisa adota abordagem quantitativa, de natureza descritiva e comparativa, com
uso de dados secundários de companhias abertas. O objetivo é verificar se empresas
de controle familiar e não familiar diferem quanto ao nível e à composição do
endividamento, e não estimar uma relação causal entre a estrutura de controle e a
estrutura de capital.

A proposta inicial previa um modelo de painel dinâmico estimado por System-GMM,
aplicado a todas as companhias não financeiras da B3. Ao longo da coleta, a
classificação manual das empresas a partir dos Formulários de Referência levou à
delimitação de um único setor (\autoref{sec:amostra}), o que resultou em uma
amostra de 16 companhias. Com esse número de unidades, estimadores de painel
dinâmico baseados em instrumentos internos não são adequados, pois exigem um
número de empresas consideravelmente maior que o de instrumentos. Optou-se, por
isso, pela comparação entre grupos por meio de estatística descritiva e de testes
de diferença entre grupos, complementados por análises de subperíodos e de
robustez.

% ---------------------------------------------------------------------
\section{Definição da amostra e fontes de dados}\label{sec:amostra}
% ---------------------------------------------------------------------

A população do estudo compreende as companhias abertas classificadas pela B3 no
setor econômico de Saúde. A restrição a um único setor tem duas finalidades:
reduzir a influência de diferenças setoriais na estrutura de capital, que a
literatura aponta como determinante relevante do endividamento
\cite{oliveira2020}, e viabilizar a classificação manual de cada companhia quanto
ao controle familiar, que exige a leitura individual dos Formulários de
Referência.
\verificar{incluir outro motivo para a escolha do setor de saúde, se houver}

Os dados contábeis foram extraídos da base Economática, em valores consolidados,
com frequência trimestral, de 31/12/2018 a 31/12/2025 (29 datas-base). Os
indicadores de fluxo, como EBITDA e despesa financeira, correspondem aos últimos
12 meses de cada data-base. Das 22 companhias do setor com dados na base, foram
mantidas apenas as que possuem informações em todo o intervalo, para que os dois
grupos sejam comparados no mesmo período. Foram excluídas seis empresas que
abriram capital ou passaram a ter dados durante o período: D1000 (DMVF, dados a
partir de 12/2019), Eurofarma (EUFA), Oncoclínicas (ONCO) e Viveo (VVEO), a partir
de 12/2020, Mater Dei (MATD), a partir de 09/2021, e Pague Menos (PGMN), a partir
de 09/2022. A amostra final é composta por 16 companhias, distribuídas em quatro
subsetores da B3 (\autoref{qua:amostra}).

\begin{quadro}[htb]
\caption{Composição da amostra por subsetor e classificação (Cenário 3)}\label{qua:amostra}
\centering\small
\begin{tabularx}{\textwidth}{|P{4.6cm}|L|L|}
\hline
\textbf{Subsetor (B3)} & \textbf{Familiares} & \textbf{Não familiares} \\ \hline
Comércio e Distribuição & Dimed/Panvel (PNVL), Profarma (PFRM), RaiaDrogasil (RADL) & -- \\ \hline
Medicamentos e Outros Produtos & Blau (BLAU), Hypera (HYPE), Ourofino (OFSA) & Biomm (BIOM) \\ \hline
Serviços Médico-Hospitalares, Análises e Diagnósticos & Dasa (DASA), Hapvida (HAPV), Rede D'Or (RDOR) & Alliar (AALR), Bradsaúde (SAUD), Fleury (FLRY), Qualicorp (QUAL) \\ \hline
Equipamentos & Baumer (BALM) & Lifemed (LMED) \\ \hline
\textbf{Total} & \textbf{10} & \textbf{6} \\ \hline
\end{tabularx}
\fonte{Elaborado pelo autor (2026), com base na classificação setorial da B3 e nos
Formulários de Referência (CVM).}
\end{quadro}

A série da Bradsaúde (SAUD) corresponde, no período analisado, à antiga Odontoprev
(ODPV3), cuja denominação e código de negociação foram alterados. A companhia não
apresenta dívida financeira em nenhuma das datas-base e foi mantida na amostra
principal, com análise de robustez que a exclui (\autoref{sec:procedimentos}).
\verificar{confirmar na DFP a ausência de empréstimos e financiamentos da SAUD e a
dívida bruta nula da HAPV em 31/12/2018}

% ---------------------------------------------------------------------
\section{Classificação das empresas familiares}\label{sec:classificacao}
% ---------------------------------------------------------------------

Não há definição única de empresa familiar na literatura, e os estudos brasileiros
combinam, de formas diferentes, critérios de propriedade, de controle e de
participação da família na administração. Para esta pesquisa, foram reunidos os
critérios de seis trabalhos (\autoref{qua:estudos}), e a partir deles foram
construídos três cenários de classificação, que diferem principalmente no
percentual mínimo de ações ordinárias detido pela família.

\begin{quadro}[htb]
\caption{Critérios de classificação de empresa familiar na literatura}\label{qua:estudos}
\centering\small
\begin{tabularx}{\textwidth}{|P{3.4cm}|L|L|P{3.2cm}|}
\hline
\textbf{Estudo} & \textbf{Participação mínima} & \textbf{Envolvimento na administração} & \textbf{Amostra} \\ \hline
\citeonline{santos2018} & Mais de 15\% das ações ordinárias, com a família no controle & Familiar no conselho de administração & 110 empresas da B3, 2010--2016 \\ \hline
\citeonline{tommasetti2019} & Mais de 50\% dos votos & Familiar como diretor-presidente ou presidente do conselho & 151 empresas da B3, 2012--2016 \\ \hline
\citeonline{voelcker2020} & Família como acionista controladora declarada & Membro indicado pelo controlador ao conselho & 1.188 observações da B3, 2010--2017 \\ \hline
\citeonline{poletti2022} & Pelo menos 20\% das ações & Familiar conselheiro, diretor-presidente ou presidente do conselho & 187 empresas da América Latina, 2005--2016 \\ \hline
\citeonline{ribeiro2026} & Pelo menos 10\% das ações ordinárias & Familiar na administração (dimensões tratadas separadamente) & 105 empresas da B3, 2010--2024 \\ \hline
\citeonline{ey2024} & Pelo menos 32\% dos votos (companhia aberta) & Não exige & 500 maiores empresas familiares do mundo \\ \hline
\end{tabularx}
\fonte{Elaborado pelo autor (2026), com base nos estudos citados.}
\end{quadro}

Em todos os cenários, a empresa é classificada como familiar somente se atender,
simultaneamente, a três condições: (i) a família detém um percentual mínimo das
ações ordinárias, somando a participação direta e a indireta; (ii) a família
integra o controle declarado da companhia; e (iii) há membro da família na
administração. A empresa que deixa de atender a pelo menos uma condição é
classificada como não familiar. O \autoref{qua:cenarios} apresenta os três
cenários.

\begin{quadro}[htb]
\caption{Cenários de classificação adotados}\label{qua:cenarios}
\centering\small
\begin{tabularx}{\textwidth}{|P{2.2cm}|P{2.2cm}|L|L|P{2.3cm}|}
\hline
\textbf{Cenário} & \textbf{Ações ordinárias da família} & \textbf{Administração} & \textbf{Base} & \textbf{Familiares $\times$ não familiares} \\ \hline
Cenário 1 & Mais de 50\% & Familiar como diretor-presidente ou presidente do conselho & \citeonline{tommasetti2019} & 5 $\times$ 11 \\ \hline
Cenário 2 & Pelo menos 32\% & Familiar no conselho de administração & \citeonline{ey2024}; \citeonline{voelcker2020} & 7 $\times$ 9 \\ \hline
\textbf{Cenário 3 (principal)} & \textbf{Mais de 15\%} & \textbf{Familiar no conselho de administração} & \textbf{\citeonline{santos2018}} & \textbf{10 $\times$ 6} \\ \hline
\end{tabularx}
\fonte{Elaborado pelo autor (2026). Em todos os cenários, exige-se também que a
família integre o controle declarado da companhia.}
\end{quadro}

O Cenário 3, que reproduz o critério de \citeonline{santos2018}, foi adotado como
classificação principal por três razões de método: trata-se de critério de estudo
brasileiro que utiliza a mesma fonte de informação desta pesquisa (Formulário de
Referência) e a mesma base de participação (ações ordinárias); capta o controle
exercido por blocos familiares com participação inferior a 50\%, frequentemente
organizados por acordo de acionistas, situação comum entre as companhias da B3; e,
de forma secundária, produz grupos de tamanho mais equilibrado, o que é relevante em uma amostra
pequena. Os Cenários 1 e 2 são utilizados como testes de robustez, para verificar
se os resultados se mantêm quando se exige uma participação familiar maior.

Para aplicar as condições, adotaram-se as seguintes regras, comuns aos três
cenários:

\begin{alineas}
\item \textbf{base de participação:} considera-se o percentual das ações
ordinárias, isto é, das ações com direito a voto, e não do capital total;

\item \textbf{grupo familiar:} a ``família'' pode corresponder a uma ou mais
famílias que integram o mesmo bloco de controle, inclusive por acordo de
acionistas, nos termos do art.~116 da Lei n.º 6.404/1976 \cite{brasil1976}. É o
caso da RaiaDrogasil (RADL), controlada por um bloco formado por três famílias
que, em conjunto, detêm cerca de 28\% das ações ordinárias e têm membros no
conselho de administração, o que a classifica como familiar no Cenário 3, mas não
nos Cenários 1 e 2;

\item \textbf{participação indireta:} soma-se a participação detida por meio de
\textit{holdings}, fundos de investimento em participações ou outras pessoas
jurídicas controladas pela família, seguindo a cadeia societária até as pessoas
físicas. Na Dasa (DASA), por exemplo, a participação da família Bueno, de cerca de
57\% das ações ordinárias, é detida por meio de um fundo de investimento em
participações;

\item \textbf{família no controle:} a participação relevante não basta se a
família não integra o controle declarado. Na Fleury (FLRY), a família Pardini
detém cerca de 20\% das ações ordinárias desde a incorporação do Instituto Hermes
Pardini, mas não faz parte do bloco de controle, e por isso a companhia é
classificada como não familiar em todos os cenários. Também são não familiares as
companhias de capital disperso ou controladas por instituições, como a Bradsaúde
(SAUD), controlada pelo Banco Bradesco;

\item \textbf{fonte e data de referência:} as informações foram obtidas no
Formulário de Referência mais recente disponível no site da Comissão de Valores
Mobiliários (CVM), nos itens de controle e composição acionária (itens 15.1 e 15.2
até 2022 e itens 6.1 e 6.2 a partir de 2023) e de administração, complementadas
pelos sites de relações com investidores. A classificação é mantida fixa em todo o
período, como em \citeonline{tommasetti2019}. Mudanças relevantes de controle ao
longo do período, como as reorganizações de Hapvida, Rede D'Or e Fleury em
2021--2022, são reconhecidas como limitação.
\end{alineas}

\verificar{percentuais de RADL (28,3\%), DASA (56,7\%) e FLRY (19,73\%) no
Formulário de Referência mais recente}

% ---------------------------------------------------------------------
\section{Indicadores de endividamento}\label{sec:indicadores}
% ---------------------------------------------------------------------

Na ausência de dados de valor de mercado e de um modelo de regressão, o
endividamento foi medido por indicadores contábeis calculados pela Economática,
organizados em três níveis de prioridade conforme a sua relação com a hipótese
(\autoref{qua:indicadores}).

\begin{quadro}[htb]
\caption{Indicadores de endividamento}\label{qua:indicadores}
\centering\small
\begin{tabularx}{\textwidth}{|P{2.5cm}|P{2.9cm}|L|L|}
\hline
\textbf{Prioridade} & \textbf{Indicador} & \textbf{Cálculo} & \textbf{O que mede} \\ \hline
Principal & Dívida bruta / Ativo total & (Empréstimos, financiamentos e debêntures de curto e longo prazo) $\div$ Ativo total & Endividamento oneroso: parcela do ativo financiada por credores financeiros, que impõem \textit{covenants} e monitoramento \\ \hline
Principal & Exigível total / Ativo total & (Passivo circulante + Passivo não circulante) $\div$ Ativo total & Endividamento geral, incluindo obrigações operacionais (fornecedores, salários, tributos, provisões) \\ \hline
Complementar & Dívida bruta / (Dívida bruta + PL) & Dívida bruta $\div$ (Dívida bruta + Patrimônio líquido) & Estrutura de capital: proporção de dívida no capital investido \\ \hline
Complementar & Dívida líquida / EBITDA & (Dívida bruta $-$ Caixa e aplicações) $\div$ EBITDA dos últimos 12 meses & Capacidade de pagamento; métrica usual em \textit{covenants} financeiros \\ \hline
Complementar & Dívida de curto prazo / Dívida total & Dívida onerosa de curto prazo $\div$ Dívida onerosa total & Composição da dívida por prazo e exposição ao risco de refinanciamento \\ \hline
Secundário & Exigível total / PL; Dívida bruta / PL; Dívida líquida / PL & Passivo ou dívida $\div$ Patrimônio líquido & Grau de endividamento em relação ao capital próprio \\ \hline
\end{tabularx}
\fonte{Elaborado pelo autor (2026), com base nas definições da Economática.}
\end{quadro}

A Dívida bruta / Ativo total é o indicador mais diretamente ligado à hipótese, pois
mede a dívida financeira, que traz os credores e as cláusulas restritivas
discutidos na \autoref{sec:risco}. O Exigível total / Ativo total permite verificar
se uma eventual diferença se restringe à dívida financeira ou alcança o passivo
como um todo. A composição da dívida por prazo foi incluída porque a dívida de
curto prazo exige renovações frequentes com os credores e, portanto, aumenta a
exposição da empresa à sua interferência \cite{diamond1991}.

Alguns indicadores não podem ser calculados em determinadas situações, e esses
casos foram tratados da seguinte forma: a Dívida de curto prazo / Dívida total não
é definida para a Bradsaúde, que não possui dívida financeira, e essa empresa é
excluída apenas desse indicador; a Dívida líquida / EBITDA não é informada pela
base quando o EBITDA é negativo, o que ocorre com a Biomm de 2018 a 2024. Nesse
indicador, a média da Biomm corresponde apenas ao exercício de 2025, e os resultados
são apresentados com e sem a empresa, acompanhados das medianas dos grupos, que são
menos sensíveis a valores extremos.
Os indicadores calculados sobre o patrimônio líquido (PL) são tratados como
secundários, por se tornarem instáveis quando o PL é pequeno ou negativo, como
ocorre com a Alliar em 2025.

% ---------------------------------------------------------------------
\section{Procedimentos de análise}\label{sec:procedimentos}
% ---------------------------------------------------------------------

A análise foi realizada em cinco etapas. Os cálculos foram feitos no Google
Planilhas e repetidos, de forma independente, em um \textit{script} em Python que
lê a mesma base e gera as tabelas deste trabalho; os dois procedimentos produzem
os mesmos resultados. A planilha de apoio, com a base de dados, a classificação
das empresas e as fórmulas, e o \textit{script} acompanham este trabalho.
\verificar{forma de disponibilização da planilha à banca: anexo, apêndice ou link}

\textbf{a) Análise descritiva.} Para cada indicador, calculou-se a média de cada
grupo em cada data-base, como média simples entre as empresas, de modo que todas as
companhias tenham o mesmo peso. A evolução dos grupos é apresentada em gráficos com
as 29 datas trimestrais, e as tabelas registram os valores de 31 de dezembro de
cada exercício, bem como o número de exercícios em que a média das familiares
ficou abaixo da média das não familiares.

\textbf{b) Testes de diferença entre os grupos.} A unidade de análise dos testes é
a empresa: para cada companhia, calculou-se a média do indicador nas oito datas de
encerramento de exercício (31/12/2018 a 31/12/2025), obtendo-se uma observação por
empresa. Esse procedimento foi adotado porque as observações de uma mesma empresa
em trimestres ou anos consecutivos não são independentes entre si, e tratá-las
como observações distintas superestimaria a significância das diferenças. Sobre
essas médias, aplicaram-se dois testes bicaudais:

\begin{itemize}
\item o teste \textit{t} de Welch \cite{welch1947}, que compara as médias dos dois
grupos sem exigir variâncias iguais nem grupos de mesmo tamanho; e
\item o teste U de Mann-Whitney \cite{mann1947}, não paramétrico, que compara as
posições (postos) das empresas de cada grupo e é menos sensível a valores
extremos e à ausência de normalidade, comuns em indicadores contábeis de amostras
pequenas. O p-valor foi obtido pela aproximação normal com correção de
continuidade; o p-valor exato, calculado para conferência, não altera nenhuma das
conclusões.
\end{itemize}

Adotaram-se como referência os níveis de significância de 1\%, 5\% e 10\%, e um
p-valor entre 5\% e 10\% é interpretado como evidência fraca. Embora a hipótese H1
seja direcional, os testes bicaudais foram mantidos por serem mais conservadores e
por a literatura apresentar argumentos nos dois sentidos (\autoref{sec:risco}).
Dado o número reduzido de empresas, que limita o poder dos testes, os resultados
não significativos são interpretados como ausência de evidência estatística de
diferença, e não como evidência de que os grupos são iguais
\cite{wasserstein2016}.

\textbf{c) Análise por subperíodos.} Para examinar se a diferença entre os grupos
se altera diante de um choque externo, o período foi dividido em três
subperíodos: pré-pandemia (2018--2019), pandemia de COVID-19 (2020--2021) e
pós-pandemia (2022--2025). A pandemia exigiu decisões de financiamento rápidas e
fora do planejamento das empresas, o que a torna um período de interesse para
observar a relação das empresas com os credores. Os testes da etapa (b) foram
repetidos com a média de cada empresa em cada subperíodo. Essa divisão foi definida
durante a análise, depois de observados os dados, e tem, por isso, caráter
exploratório: os resultados dos três subperíodos são reportados integralmente,
inclusive os não significativos, e, para o conjunto desses testes, informa-se o
p-valor ajustado pelo método de Holm.

\textbf{d) Análises de robustez.} Os testes foram repetidos (i) com as
classificações dos Cenários 1 e 2, para verificar a sensibilidade dos resultados ao
critério de empresa familiar; (ii) excluindo a Bradsaúde, operadora de planos
controlada por uma instituição financeira e única companhia da amostra sem dívida
financeira, para avaliar a influência dessa observação sobre o resultado, sem que
isso signifique tratar a ausência de dívida como erro, já que ela é uma escolha
legítima de financiamento; (iii) retirando uma empresa de cada vez, entre as 16,
para verificar se algum resultado depende de uma única companhia; e (iv) excluindo
as três empresas do subsetor de Comércio e Distribuição, em que não há empresas
não familiares.

\textbf{e) Magnitude das diferenças.} Além dos p-valores, são apresentados a
diferença entre as médias dos grupos, o seu intervalo de confiança de 95\% e o
tamanho de efeito (g de Hedges), para indicar a magnitude e a imprecisão das
estimativas.

Como foram realizados diversos testes (vários indicadores, dois testes,
subperíodos e cenários), um resultado significativo isolado deve ser interpretado
com cautela, e as conclusões consideram o conjunto dos resultados, e não testes
individuais.

% [OPCIONAL – só incluir se a regressão for refeita com o subsetor corrigido e um
% script reprodutível (R):]
% \textbf{e) Análise exploratória multivariada.} Como complemento, estimou-se uma
% regressão linear com dados empilhados (\textit{pooled}), tendo como variável
% dependente a Dívida bruta / Ativo total e como variável de interesse uma variável
% binária para empresa familiar, controlando por tamanho (logaritmo do capital
% empregado), rentabilidade do ativo, liquidez corrente, imobilização do patrimônio
% líquido e efeitos de ano, com erros-padrão agrupados por empresa. Dado o número
% de empresas, os resultados dessa etapa são tratados como exploratórios.

% =====================================================================
\chapter{Resultados e discussão}\label{cap:resultados}
% =====================================================================
% Todos os números deste capítulo saem de tcc/analise/resultados_c3.py, que lê a base
% da Economática e gera as tabelas em tabelas/*.tex. Se a base mudar, rode o script de
% novo e confira os números citados no texto (procure por "%" seguido de "fonte:").

Este capítulo apresenta a evolução do endividamento dos dois grupos, os testes de
diferença no período completo e nos subperíodos e as análises de robustez, e
discute os resultados à luz do referencial teórico.

% ---------------------------------------------------------------------
\section{Evolução do endividamento dos grupos}\label{sec:evolucao}
% ---------------------------------------------------------------------

A \autoref{tab:anuais} apresenta as médias anuais dos três indicadores centrais da
análise. No endividamento oneroso, medido pela Dívida bruta / Ativo total, as
empresas familiares ficaram abaixo das não familiares em todos os oito exercícios.
A distância, porém, não foi constante: em 2018 e 2019, os grupos estavam próximos
(23,9\% contra 26,4\% e 27,0\% contra 28,3\%); a diferença se ampliou em 2021 e 2022,
quando as não familiares chegaram a 35,4\% e 36,1\%, contra 26,0\% e 24,9\% das
familiares; e voltou a se estreitar a partir de 2024, até quase desaparecer em 2025
(26,4\% contra 26,9\%).
% fonte: saidas/resultados.md, seção 2

% Gerado por tcc/analise/resultados_c3.py -- não edite à mão; rode o script de novo.
\begin{table}[htb]
\caption{Médias anuais por grupo em 31 de dezembro (\%)}\label{tab:anuais}
\centering\footnotesize\setlength{\tabcolsep}{4pt}
\begin{tabular}{lrrrrrr}
\toprule
Ano & \multicolumn{2}{c}{Dívida bruta / Ativo} & \multicolumn{2}{c}{Exigível / Ativo} & \multicolumn{2}{c}{Dívida CP / Dívida total} \\
\midrule
 & F & NF & F & NF & F & NF \\
2018 & 23,9 & 26,4 & 50,2 & 47,8 & 29,2 & 41,4 \\
2019 & 27,0 & 28,3 & 56,1 & 53,3 & 23,0 & 25,4 \\
2020 & 29,6 & 30,5 & 59,0 & 56,0 & 17,2 & 25,8 \\
2021 & 26,0 & 35,4 & 54,5 & 60,1 & 21,9 & 47,0 \\
2022 & 24,9 & 36,1 & 54,9 & 59,0 & 19,6 & 35,9 \\
2023 & 26,2 & 32,8 & 55,5 & 56,8 & 18,9 & 36,8 \\
2024 & 24,7 & 28,2 & 55,6 & 54,6 & 16,1 & 39,1 \\
2025 & 26,4 & 26,9 & 57,2 & 62,1 & 19,1 & 54,3 \\
\bottomrule
\end{tabular}
\fonte{Elaborado pelo autor (2026), com dados da Economática.}
\nota{F: familiares (10 empresas); NF: não familiares (6 empresas; 5 na Dívida CP / Dívida total). Média simples entre as empresas de cada grupo.}
\end{table}


O endividamento geral, medido pelo Exigível total / Ativo total, não seguiu o
mesmo padrão. As familiares tiveram médias superiores às das não familiares em
quatro dos oito exercícios, incluindo os anos anteriores à pandemia, e as médias
dos dois grupos ficaram entre 47,8\% e 62,1\% ao longo de todo o período. A
diferença observada na dívida onerosa, portanto, não corresponde a um passivo
total menor, e sim a uma composição do passivo com menos dívida financeira.

A composição da dívida por prazo apresentou a diferença mais persistente. Em todos
os exercícios, a participação da dívida de curto prazo foi menor nas familiares, e
a distância se tornou expressiva a partir de 2021: desde então, a dívida de curto
prazo representou entre 16\% e 22\% da dívida total das familiares, contra 36\% a
54\% nas não familiares.

% Gráficos do autor (exportar do Google Planilhas para figuras/ como PDF ou PNG):
% \begin{figure}[htb]
%   \caption{Dívida bruta / Ativo total: média trimestral por grupo (2018--2025)}\label{fig:divbruta}
%   \centering
%   \includegraphics[width=\textwidth]{figuras/div-bruta-ativo.pdf}
%   \fonte{Elaborado pelo autor (2026), com dados da Economática.}
% \end{figure}
% (repetir para Exigível / Ativo e Dívida CP / Dívida total)
\verificar{inserir os gráficos trimestrais dos três indicadores (Figuras 1 a 3) e
citá-los neste ponto do texto}

% ---------------------------------------------------------------------
\section{Comparação entre os grupos no período completo}\label{sec:completo}
% ---------------------------------------------------------------------

A \autoref{tab:completo} apresenta os testes de diferença para o período de 2018 a
2025. Nos cinco indicadores, a média das familiares foi inferior à das não
familiares, mas nenhuma diferença foi significativa ao nível de 5\% em qualquer dos
dois testes.

% Gerado por tcc/analise/resultados_c3.py -- não edite à mão; rode o script de novo.
\begin{table}[htb]
\caption{Endividamento das empresas familiares (F) e não familiares (NF), 2018--2025}\label{tab:completo}
\centering\footnotesize\setlength{\tabcolsep}{4pt}
\begin{tabular}{lrrcccc}
\toprule
Indicador & F & NF & n & Anos F $<$ NF & p (t) & p (MW) \\
\midrule
Dívida bruta / Ativo total & 26,1\% & 30,6\% & 10 $\times$ 6 & 8 de 8 & 0,556 & 0,255 \\
Exigível total / Ativo total & 55,4\% & 56,2\% & 10 $\times$ 6 & 4 de 8 & 0,885 & 0,957 \\
Dívida bruta / (Dívida bruta + PL) & 37,2\% & 43,5\% & 10 $\times$ 6 & 6 de 8 & 0,563 & 0,255 \\
Dívida líquida / EBITDA & 1,4x & 4,4x & 10 $\times$ 6 & 5 de 8 & 0,285 & 0,416 \\
Dívida CP / Dívida total & 20,4\% & 38,2\% & 10 $\times$ 5 & 8 de 8 & 0,112 & 0,058$^{*}$ \\
\bottomrule
\end{tabular}
\fonte{Elaborado pelo autor (2026), com dados da Economática.}
\nota{Testes bicaudais sobre a média de cada empresa nas datas de 31/12 do período. p (MW): teste de Mann-Whitney, aproximação normal com correção de continuidade. $^{***}$, $^{**}$ e $^{*}$: significativo a 1\%, 5\% e 10\%. Valores: média simples das médias por empresa. Na Dívida CP / Dívida total, a Bradsaúde é excluída por não ter dívida financeira; na Dívida líquida / EBITDA, a Biomm tem um único ano com EBITDA positivo (2025).}
\end{table}


No indicador principal, a Dívida bruta / Ativo total, a média das familiares foi
de 26,1\%, contra 30,6\% das não familiares, diferença de 4,5 pontos percentuais que
não foi estatisticamente significativa (p = 0,556 no teste t e p = 0,255 no teste
de Mann-Whitney). O Exigível total / Ativo total praticamente não diferiu entre os
grupos (55,4\% e 56,2\%; p = 0,885 e 0,957). A estrutura de capital, medida pela
Dívida bruta / (Dívida bruta + PL), seguiu a mesma direção da dívida onerosa
(37,2\% contra 43,5\%), também sem significância.

O resultado mais próximo da significância foi o da composição da dívida: a dívida
de curto prazo representou, em média, 20,4\% da dívida total das familiares e
38,2\% da dívida das não familiares. A diferença foi significativa apenas ao nível
de 10\% no teste de Mann-Whitney (p = 0,058) e não foi significativa no teste t
(p = 0,112). A divergência entre os testes é coerente com a dispersão do grupo não
familiar, em que a Lifemed concentrou 66\% de sua dívida no curto prazo, valor que
eleva a variância e reduz o poder do teste t.
\verificar{conferir no Google Planilhas o valor da LMED (média de 66,3\%) antes de
citar}

Na Dívida líquida / EBITDA, a média das não familiares (4,4x) foi três vezes a das
familiares (1,4x), mas esse contraste se deve, em grande parte, à Biomm, cujo
indicador só pôde ser calculado em 2025 (16,1x), já que o EBITDA da empresa foi
negativo nos anos anteriores. Sem a Biomm, a média das não familiares cai para
2,0x (p = 0,560 no teste t), e as medianas dos grupos são 1,6x e 2,6x. Esse
indicador, portanto, não permite concluir que as familiares tenham maior
capacidade de pagamento.
\decidir{confirmar este tratamento da Biomm (apresentar o valor com e sem a empresa
e as medianas) e ajustar a seção 3.4 de acordo}

% ---------------------------------------------------------------------
\section{Subperíodos: antes, durante e depois da pandemia}\label{sec:subperiodos}
% ---------------------------------------------------------------------

A \autoref{tab:subperiodos} repete os testes em cada subperíodo. Antes da pandemia
(2018--2019), os grupos foram semelhantes em todos os indicadores, e nenhuma
diferença se aproximou da significância.

% Gerado por tcc/analise/resultados_c3.py -- não edite à mão; rode o script de novo.
\begin{table}[htb]
\caption{Comparação entre os grupos por subperíodo}\label{tab:subperiodos}
\centering\footnotesize\setlength{\tabcolsep}{4pt}
\begin{tabular}{llrrcc}
\toprule
Indicador & Subperíodo & F & NF & p (t) & p (MW) \\
\midrule
Dív. bruta / Ativo & 2018--2019 & 25,4\% & 27,3\% & 0,815 & 0,625 \\
 & 2020--2021 & 27,8\% & 33,0\% & 0,543 & 0,255 \\
 & 2022--2025 & 25,5\% & 31,0\% & 0,497 & 0,416 \\
Exigível / Ativo & 2018--2019 & 53,1\% & 50,6\% & 0,703 & 0,704 \\
 & 2020--2021 & 56,8\% & 58,1\% & 0,825 & 0,871 \\
 & 2022--2025 & 55,8\% & 58,1\% & 0,760 & 0,957 \\
Dív. CP / Dív. total & 2018--2019 & 24,9\% & 33,4\% & 0,435 & 0,298 \\
 & 2020--2021 & 19,5\% & 36,4\% & \textbf{0,040}$^{**}$ & \textbf{0,043}$^{**}$ \\
 & 2022--2025 & 18,4\% & 41,5\% & 0,130 & 0,098$^{*}$ \\
\bottomrule
\end{tabular}
\fonte{Elaborado pelo autor (2026), com dados da Economática.}
\nota{Testes bicaudais sobre a média de cada empresa nas datas de 31/12 do período. p (MW): teste de Mann-Whitney, aproximação normal com correção de continuidade. $^{***}$, $^{**}$ e $^{*}$: significativo a 1\%, 5\% e 10\%. Subperíodos: pré-pandemia (2018--2019), pandemia (2020--2021) e pós-pandemia (2022--2025).}
\end{table}


Durante a pandemia (2020--2021), a diferença na composição da dívida tornou-se
significativa ao nível de 5\% nos dois testes: a dívida de curto prazo representou
19,5\% da dívida total das familiares e 36,4\% da das não familiares (p = 0,040 no
teste t e p = 0,043 no teste de Mann-Whitney). No mesmo período, a dívida onerosa
das familiares também foi menor (27,8\% contra 33,0\% do ativo), mas sem
significância (p = 0,543 e 0,255), e o endividamento geral permaneceu semelhante
(56,8\% e 58,1\%).

Após a pandemia (2022--2025), a distância na composição da dívida aumentou em
termos absolutos (18,4\% contra 41,5\%), mas a significância se reduziu: a
diferença foi significativa apenas a 10\% no teste de Mann-Whitney (p = 0,098) e
não foi significativa no teste t (p = 0,130).

Como explicado na \autoref{sec:procedimentos}, a divisão em subperíodos tem caráter
exploratório, e o resultado da pandemia é o único significativo a 5\% entre os
testes realizados. Ele deve, por isso, ser lido como um indício compatível com a
hipótese, e não como sua confirmação.

% ---------------------------------------------------------------------
\section{Análises de robustez}\label{sec:robustez}
% ---------------------------------------------------------------------

A \autoref{tab:robustez} mostra como os resultados se alteram com a exclusão da
Bradsaúde e com as classificações alternativas de empresa familiar.

% Gerado por tcc/analise/resultados_c3.py -- não edite à mão; rode o script de novo.
\begin{table}[htb]
\caption{Análises de robustez}\label{tab:robustez}
\centering\footnotesize\setlength{\tabcolsep}{4pt}
\begin{tabular}{llrrccc}
\toprule
Especificação & Indicador & F & NF & n & p (t) & p (MW) \\
\midrule
C3 (principal) & Dív. bruta / Ativo & 26,1\% & 30,6\% & 10 $\times$ 6 & 0,556 & 0,255 \\
 & Exigível / Ativo & 55,4\% & 56,2\% & 10 $\times$ 6 & 0,885 & 0,957 \\
 & Dív. CP / Dív. total & 20,4\% & 38,2\% & 10 $\times$ 5 & 0,112 & 0,058$^{*}$ \\
C3 sem Bradsaúde & Dív. bruta / Ativo & 26,1\% & 36,7\% & 10 $\times$ 5 & \textbf{0,027}$^{**}$ & 0,058$^{*}$ \\
 & Exigível / Ativo & 55,4\% & 59,8\% & 10 $\times$ 5 & 0,348 & 0,759 \\
 & Dív. CP / Dív. total & 20,4\% & 38,2\% & 10 $\times$ 5 & 0,112 & 0,058$^{*}$ \\
C2 & Dív. bruta / Ativo & 28,7\% & 27,0\% & 7 $\times$ 9 & 0,801 & 1,000 \\
 & Exigível / Ativo & 53,6\% & 57,3\% & 7 $\times$ 9 & 0,567 & 0,597 \\
 & Dív. CP / Dív. total & 19,1\% & 32,6\% & 7 $\times$ 8 & 0,080$^{*}$ & 0,118 \\
C1 & Dív. bruta / Ativo & 30,9\% & 26,3\% & 5 $\times$ 11 & 0,558 & 0,734 \\
 & Exigível / Ativo & 55,8\% & 55,7\% & 5 $\times$ 11 & 0,987 & 0,910 \\
 & Dív. CP / Dív. total & 19,2\% & 29,9\% & 5 $\times$ 10 & 0,131 & 0,245 \\
\bottomrule
\end{tabular}
\fonte{Elaborado pelo autor (2026), com dados da Economática.}
\nota{Testes bicaudais sobre a média de cada empresa nas datas de 31/12 do período. p (MW): teste de Mann-Whitney, aproximação normal com correção de continuidade. $^{***}$, $^{**}$ e $^{*}$: significativo a 1\%, 5\% e 10\%. C1, C2 e C3: Cenários 1, 2 e 3; C1 e C2 são classificações alternativas de empresa familiar (\autoref{qua:cenarios}).}
\end{table}


A exclusão da Bradsaúde, única empresa da amostra sem dívida financeira, eleva a
média de endividamento oneroso das não familiares de 30,6\% para 36,7\% e torna a
diferença significativa a 5\% no teste t (p = 0,027) e a 10\% no teste de
Mann-Whitney (p = 0,058). O resultado do indicador principal depende, portanto, de
uma única observação: com ela, a diferença não é significativa; sem ela, passa a
ser.

As classificações alternativas produzem um quadro distinto. Nos Cenários 1 e 2, que
exigem participação familiar maior, a diferença na dívida onerosa inverte o sinal:
as familiares passam a ter média superior à das não familiares (30,9\% contra
26,4\% no Cenário 1 e 28,7\% contra 27,0\% no Cenário 2), sem significância. A
inversão ocorre porque RaiaDrogasil, Profarma e Hypera, familiares apenas no
Cenário 3, passam ao grupo não familiar. As duas primeiras, que atuam no comércio e
na distribuição de medicamentos, estão entre as empresas menos endividadas da
amostra (12,8\% e 16,5\% do ativo, em média). Parte da diferença observada no
Cenário 3 pode refletir, portanto, o modelo de negócio desse subsetor, e não
apenas a estrutura de controle.

A composição da dívida, por outro lado, manteve a direção nos três cenários: a
participação da dívida de curto prazo foi menor nas familiares em todos eles (19,1\%
contra 32,6\% no Cenário 2, com p = 0,080 no teste t, e 19,2\% contra 29,9\% no
Cenário 1, sem significância). É, entre os indicadores analisados, o resultado
menos sensível ao critério de classificação.

% ---------------------------------------------------------------------
\section{Discussão}\label{sec:discussao}
% ---------------------------------------------------------------------

Os resultados não permitem confirmar a hipótese H1. Embora as empresas familiares
tenham apresentado menor endividamento oneroso em todos os exercícios, a diferença
não foi estatisticamente significativa na amostra principal, não se manteve sob as
classificações alternativas e dependeu da presença de uma única empresa sem dívida.
Não foram encontradas, portanto, evidências estatisticamente significativas de que
as empresas familiares do setor de saúde da B3 se endividem menos que as não
familiares.

Esse resultado é coerente com a literatura inconclusiva apresentada no
\autoref{cap:referencial}. As duas forças discutidas na \autoref{sec:risco} podem
atuar em sentidos opostos: a aversão à interferência de credores, associada à
preservação da riqueza socioemocional \cite{santos2018,ames2020}, levaria a menos
dívida, enquanto o receio de diluir o controle levaria a preferir a dívida à emissão
de ações \cite{poletti2022}. No Brasil, \citeonline{kayo2018} encontraram empresas
familiares mais endividadas, o que indica que a segunda força pode ser relevante no
mercado brasileiro. Em uma amostra pequena, a combinação dessas forças é compatível
com diferenças de pequena magnitude e instáveis entre especificações, como as
observadas.

O resultado mais consistente refere-se ao prazo, e não ao nível, da dívida. As
familiares concentraram menos dívida no curto prazo em todos os exercícios e nos
três critérios de classificação, e a diferença foi significativa durante a
pandemia. Esse padrão é compatível com o argumento de \citeonline{diamond1991}: a
dívida de curto prazo precisa ser renovada com frequência e expõe a empresa ao
risco de que os credores não a refinanciem, risco que se agrava em períodos de
crise. Para uma família que busca preservar a autonomia decisória, alongar a dívida
reduz a necessidade de renegociações frequentes e, com ela, a exposição à avaliação
dos credores. Assim, a aversão à interferência de credores parece se manifestar
mais na escolha do prazo do que na quantidade de dívida.

Dois cuidados limitam essa interpretação. Primeiro, a diferença de prazo foi
significativa a 5\% apenas em um subperíodo definido durante a análise, entre
diversos testes realizados, o que aumenta a chance de um resultado significativo
ocorrer por acaso. Segundo, o desenho descritivo não permite separar o efeito do
controle familiar de outras características das empresas, como o subsetor de
atuação \cite{oliveira2020} e o tamanho. Os resultados devem, por isso, ser lidos como evidência
descritiva de uma postura mais conservadora das famílias em relação ao prazo da
dívida, a ser testada em amostras maiores e com controles.

% =====================================================================
\chapter{Considerações finais}\label{cap:consideracoes}
% =====================================================================
% A ESCREVER. Base (síntese da orientação de 05/11/2026, seção 5.1), limitações e
% pesquisas futuras: ver tcc/CONTEXTO_TCC.md, seção 12, itens 3 a 5.

\textcolor{red}{[Capítulo a escrever.]}


% ---------------------------------------------------------------------
% Elementos pós-textuais
% ---------------------------------------------------------------------
\postextual
\bibliography{referencias}

\begin{apendicesenv}
\partapendices
% =====================================================================
\chapter{Classificação das empresas da amostra}\label{ap:classificacao}
% =====================================================================

O \autoref{qua:porempresa} apresenta, para cada empresa da amostra, a participação
aproximada da família (ou do grupo de famílias) nas ações ordinárias, somadas a
participação direta e a indireta, e a classificação resultante em cada cenário.
\verificar{conferir cada percentual, o item e a data do Formulário de Referência
consultado, e completar o tipo de controlador das não familiares}

\begin{quadro}[htb]
\caption{Participação familiar e classificação por cenário}\label{qua:porempresa}
\centering\footnotesize
\begin{tabularx}{\textwidth}{|l|P{2.3cm}|c|c|c|L|}
\hline
\textbf{Código} & \textbf{Ações ordinárias da família (aprox.)} & \textbf{C1} & \textbf{C2} & \textbf{C3} & \textbf{Observação} \\ \hline
BLAU & 82,5\% & F & F & F & Abertura de capital em 2021 \\ \hline
BALM & 89\% & F & F & F & Familiar preside o conselho e a diretoria \\ \hline
DASA & 56,7\% & F & F & F & Participação por fundo de investimento em participações \\ \hline
OFSA & 56,7\% & F & F & F & Bloco de duas famílias \\ \hline
RDOR & 50,2\% & F & F & F & No limite do Cenário 1 \\ \hline
HAPV & 38\% & NF & F & F & Acima de 50\% até a combinação com a NotreDame Intermédica \\ \hline
PNVL & 38,8\% & NF & F & F & Bloco de três famílias \\ \hline
RADL & 28,3\% & NF & NF & F & Bloco de três famílias \\ \hline
PFRM & 28\% & NF & NF & F & Controle compartilhado com investidor estratégico \\ \hline
HYPE & 25,9\% & NF & NF & F & \\ \hline
FLRY & 19,7\% & NF & NF & NF & Família não integra o controle \\ \hline
BIOM & 8\% & NF & NF & NF & \\ \hline
AALR & -- & NF & NF & NF & Sem família no controle \\ \hline
QUAL & -- & NF & NF & NF & Sem família no controle \\ \hline
LMED & -- & NF & NF & NF & Sem família no controle \\ \hline
SAUD & -- & NF & NF & NF & Controlada por instituição financeira \\ \hline
\end{tabularx}
\fonte{Elaborado pelo autor (2026), com base nos Formulários de Referência (CVM).}
\nota{F: familiar; NF: não familiar. C1, C2 e C3: Cenários 1, 2 e 3 (\autoref{qua:cenarios}).}
\end{quadro}

% =====================================================================
\chapter{Indicadores secundários}\label{ap:secundarios}
% =====================================================================

A \autoref{tab:secundarios} apresenta os indicadores calculados sobre o patrimônio
líquido, tratados como secundários (\autoref{sec:indicadores}), e a liquidez
corrente. Nenhuma diferença é significativa a 5\%. A liquidez corrente das
familiares foi maior em todos os exercícios, com diferença significativa a 10\%
no teste t, o que é coerente com a sua menor participação de dívida de curto prazo.

% Gerado por tcc/analise/resultados_c3.py -- não edite à mão; rode o script de novo.
\begin{table}[htb]
\caption{Indicadores secundários e liquidez corrente, 2018--2025}\label{tab:secundarios}
\centering\footnotesize\setlength{\tabcolsep}{4pt}
\begin{tabular}{lrrccc}
\toprule
Indicador & F & NF & n & p (t) & p (MW) \\
\midrule
Exigível total / PL & 153,6\% & 116,7\% & 10 $\times$ 6 & 0,349 & 0,416 \\
Dívida bruta / PL & 73,0\% & 74,7\% & 10 $\times$ 6 & 0,950 & 0,786 \\
Dívida líquida / PL & 37,8\% & 41,4\% & 10 $\times$ 6 & 0,855 & 0,625 \\
Liquidez corrente & 2,2x & 1,5x & 10 $\times$ 6 & 0,064$^{*}$ & 0,116 \\
\bottomrule
\end{tabular}
\fonte{Elaborado pelo autor (2026), com dados da Economática.}
\nota{Testes bicaudais sobre a média de cada empresa nas datas de 31/12 do período. p (MW): teste de Mann-Whitney, aproximação normal com correção de continuidade. $^{***}$, $^{**}$ e $^{*}$: significativo a 1\%, 5\% e 10\%. Os indicadores sobre o patrimônio líquido (PL) incluem a Alliar em 2025, quando o PL da empresa foi negativo.}
\end{table}


\end{apendicesenv}

\end{document}
